% GENERATED FILE. Edit canonical lessons, then run npm run generate:books.
% canonical-source-hash: fnv1a64:e59bcd0985a48809
% canonical-lessons: RU-C466-burya, RU-C466-uragan, RU-C466-poryv, RU-C466-gremet, RU-C466-sverkat, RU-R466-review-459-466

\chapter{Storms}
\label{ch:ru-storms}
\hlchaptermodality{466}

\begin{chapteropening}
\textbf{By the end of this chapter:} \emph{I can talk about storms: a storm, a hurricane, a gust, thunder rumbling and lightning flashing.}

Complete the last lesson of chapter 466: I can talk about storms: a storm, a hurricane, a gust, thunder rumbling and lightning flashing.
\end{chapteropening}

\section[\texorpdfstring{\ru{буря} (búrya)}{буря (búrya)}]{\ru{буря} (búrya) — a storm}

\label{lesson:RU-C466-burya}



\begin{quote}
Before the new one: say the Russian for a puddle, then the Russian for to pour.
\end{quote}

\subsection*{You'll want to know: \ru{буря}}
\textbf{\ru{буря}} (\emph{búrya}) — \textquotedblleft{}a storm\textquotedblright{}.

\begin{quote}
\ru{Ночью} \ru{была} \ru{сильная} \ru{буря}.
\end{quote}

\begin{quote}
\emph{Nóchyu bylá sílnaya búrya.}
\end{quote}

\begin{quote}
There was a strong storm in the night.
\end{quote}

Feminine. \emph{\ru{Буря}} is a storm of wind; \emph{\ru{гроза}} is a thunderstorm.

\begin{cousinweb}
\emph{\ru{буря}} is inherited from Proto-Slavic \emph{buřa} (EtymDB).
\end{cousinweb}

\begin{grammarlens}[title={What you've built}]
One more word for telling, explaining and describing at B1.
\end{grammarlens}

\subsection*{Your turn}
\begin{itemize}
\raggedright
  \item \emph{Say it:} \emph{búrya}
  \item \emph{Say it:} \emph{búrya}, once more
  \item \emph{Say it:} a sentence of your own with \emph{búrya}
  \item \emph{Recall:} say the Russian for a puddle, then the Russian for to pour, then say \emph{búrya} again
\end{itemize}

\begin{tcolorbox}[breakable,colback=teal!4,colframe=teal!35!black,title={Before you move on}]
What does \ru{буря} mean? (\textquotedblleft{}A storm\textquotedblright{}.) Say it once more.
\end{tcolorbox}
\section[\texorpdfstring{\ru{ураган} (uragán)}{ураган (uragán)}]{\ru{ураган} (uragán) — a hurricane}

\label{lesson:RU-C466-uragan}



\begin{quote}
Before the new one: say the Russian for to pour, then the Russian for a storm.
\end{quote}

\subsection*{You'll want to know: \ru{ураган}}
\textbf{\ru{ураган}} (\emph{uragán}) — \textquotedblleft{}a hurricane\textquotedblright{}.

\begin{quote}
\ru{Ураган} \ru{сломал} \ru{много} \ru{деревьев}.
\end{quote}

\begin{quote}
\emph{Uragán slomál mnógo derévyev.}
\end{quote}

\begin{quote}
The hurricane broke a lot of trees.
\end{quote}

Masculine.

\begin{cousinweb}
\emph{\ru{ураган}} was borrowed from French \emph{ouragan}, which came from Spanish \emph{huracán} (EtymDB).
\end{cousinweb}

\begin{grammarlens}[title={What you've built}]
One more word for telling, explaining and describing at B1.
\end{grammarlens}

\subsection*{Your turn}
\begin{itemize}
\raggedright
  \item \emph{Say it:} \emph{uragán}
  \item \emph{Say it:} \emph{uragán}, once more
  \item \emph{Say it:} a sentence of your own with \emph{uragán}
  \item \emph{Recall:} say the Russian for to pour, then the Russian for a storm, then say \emph{uragán} again
\end{itemize}

\begin{tcolorbox}[breakable,colback=teal!4,colframe=teal!35!black,title={Before you move on}]
What does \ru{ураган} mean? (\textquotedblleft{}A hurricane\textquotedblright{}.) Say it once more.
\end{tcolorbox}
\section[\texorpdfstring{\ru{порыв} (porýv)}{порыв (porýv)}]{\ru{порыв} (porýv) — a gust (of wind); an impulse}

\label{lesson:RU-C466-poryv}



\begin{quote}
Before the new one: say the Russian for a storm, then the Russian for a hurricane.
\end{quote}

\subsection*{You'll want to know: \ru{порыв}}
\textbf{\ru{порыв}} (\emph{porýv}) — \textquotedblleft{}a gust (of wind); an impulse\textquotedblright{}.

\begin{quote}
\ru{Сильный} \ru{порыв} \ru{ветра} \ru{сорвал} \ru{с} \ru{него} \ru{шляпу}.
\end{quote}

\begin{quote}
\emph{Sílnyy porýv vétra sorvál s negó shlyápu.}
\end{quote}

\begin{quote}
A strong gust of wind blew his hat off.
\end{quote}

Masculine. \emph{\ru{Порыв} \ru{ветра}} — a gust of wind.

\begin{cousinweb}
\emph{\ru{порыв}} is built from \emph{\ru{по}-} and \emph{\ru{рыв}} 'a jerk, a tearing', on the root of \emph{\ru{рвать}} 'to tear'. EtymDB lists \emph{\ru{порыв}} without a deeper source.
\end{cousinweb}

\begin{grammarlens}[title={What you've built}]
One more word for telling, explaining and describing at B1.
\end{grammarlens}

\subsection*{Your turn}
\begin{itemize}
\raggedright
  \item \emph{Say it:} \emph{porýv}
  \item \emph{Say it:} \emph{porýv}, once more
  \item \emph{Say it:} a sentence of your own with \emph{porýv}
  \item \emph{Recall:} say the Russian for a storm, then the Russian for a hurricane, then say \emph{porýv} again
\end{itemize}

\begin{tcolorbox}[breakable,colback=teal!4,colframe=teal!35!black,title={Before you move on}]
What does \ru{порыв} mean? (\textquotedblleft{}A gust (of wind); an impulse\textquotedblright{}.) Say it once more.
\end{tcolorbox}
\section[\texorpdfstring{\ru{греметь} (gremét)}{греметь (gremét)}]{\ru{греметь} (gremét) — to rumble, to thunder}

\label{lesson:RU-C466-gremet}



\begin{quote}
Before the new one: say the Russian for a hurricane, then the Russian for a gust.
\end{quote}

\subsection*{You'll want to know: \ru{греметь}}
\textbf{\ru{греметь}} (\emph{gremét}) — \textquotedblleft{}to rumble, to thunder\textquotedblright{}.

\begin{quote}
\ru{Всю} \ru{ночь} \ru{гремел} \ru{гром}.
\end{quote}

\begin{quote}
\emph{Vsyu noch gremél grom.}
\end{quote}

\begin{quote}
Thunder rumbled all night.
\end{quote}

Imperfective; perfective \emph{\ru{прогреметь}}. Present \emph{\ru{гремлю}, \ru{гремишь}}.

\begin{cousinweb}
\emph{\ru{греметь}} is inherited from Proto-Slavic \emph{grĭměti} (EtymDB), which EtymDB takes back to the PIE root \emph{g\textsuperscript{h}rem-} 'to resound, to thunder'; \emph{\ru{гром}} 'thunder' is its noun.
\end{cousinweb}

\begin{grammarlens}[title={What you've built}]
One more word for telling, explaining and describing at B1.
\end{grammarlens}

\subsection*{Your turn}
\begin{itemize}
\raggedright
  \item \emph{Say it:} \emph{gremét}
  \item \emph{Say it:} \emph{gremét}, once more
  \item \emph{Say it:} a sentence of your own with \emph{gremét}
  \item \emph{Recall:} say the Russian for a hurricane, then the Russian for a gust, then say \emph{gremét} again
\end{itemize}

\begin{tcolorbox}[breakable,colback=teal!4,colframe=teal!35!black,title={Before you move on}]
What does \ru{греметь} mean? (\textquotedblleft{}To rumble, to thunder\textquotedblright{}.) Say it once more.
\end{tcolorbox}
\section[\texorpdfstring{\ru{сверкать} (sverkát)}{сверкать (sverkát)}]{\ru{сверкать} (sverkát) — to flash, to sparkle}

\label{lesson:RU-C466-sverkat}



\begin{quote}
Before the new one: say the Russian for a gust, then the Russian for to rumble.
\end{quote}

\subsection*{You'll want to know: \ru{сверкать}}
\textbf{\ru{сверкать}} (\emph{sverkát}) — \textquotedblleft{}to flash, to sparkle\textquotedblright{}.

\begin{quote}
\ru{Над} \ru{морем} \ru{сверкают} \ru{молнии}.
\end{quote}

\begin{quote}
\emph{Nad mórem sverkáyut mólnii.}
\end{quote}

\begin{quote}
Lightning is flashing over the sea.
\end{quote}

Imperfective; perfective \emph{\ru{сверкнуть}}.

\begin{cousinweb}
The origin of \emph{\ru{сверкать}} is not settled; EtymDB lists the word without a source.
\end{cousinweb}

\begin{grammarlens}[title={What you've built}]
One more word for telling, explaining and describing at B1.
\end{grammarlens}

\subsection*{Your turn}
\begin{itemize}
\raggedright
  \item \emph{Say it:} \emph{sverkát}
  \item \emph{Say it:} \emph{sverkát}, once more
  \item \emph{Say it:} a sentence of your own with \emph{sverkát}
  \item \emph{Recall:} say the Russian for a gust, then the Russian for to rumble, then say \emph{sverkát} again
\end{itemize}

\begin{tcolorbox}[breakable,colback=teal!4,colframe=teal!35!black,title={Before you move on}]
What does \ru{сверкать} mean? (\textquotedblleft{}To flash, to sparkle\textquotedblright{}.) Say it once more.
\end{tcolorbox}
\section[(dialogue)]{Words from chapters 459–466 — a review}

\label{lesson:RU-R466-review-459-466}



\begin{quote}
Say the Russian for a hike and a path.
\end{quote}

\subsection*{Your turn}
Say these aloud:

\begin{itemize}
\raggedright
  \item a hike, a path, a rest stop, to scramble, to stumble
  \item a stream, a swamp, a spring, a pond, to babble
  \item a clearing, the edge of a forest, a thicket, moss, a cone
  \item a tent, a sleeping bag, a campfire, to light, to put out
  \item a high tide, a low tide, a bay, a seashell, to dive
  \item a forecast, a weather forecaster, precipitation, overcast, to clear up
  \item a downpour, to drizzle, to get soaked, a puddle, to pour
  \item a storm, a hurricane, a gust, to rumble, to flash
\end{itemize}

\begin{tcolorbox}[breakable,colback=teal!4,colframe=teal!35!black,title={Before you move on}]
A hike? (\textbf{\ru{поход}}.) A path? (\textbf{\ru{тропа}}.) A rest stop? (\textbf{\ru{привал}}.) To scramble? (\textbf{\ru{карабкаться}}.) To stumble? (\textbf{\ru{спотыкаться}}.) A stream? (\textbf{\ru{ручей}}.) A swamp? (\textbf{\ru{болото}}.) A spring? (\textbf{\ru{родник}}.) A pond? (\textbf{\ru{пруд}}.) To babble? (\textbf{\ru{журчать}}.) A clearing? (\textbf{\ru{поляна}}.) The edge of a forest? (\textbf{\ru{опушка}}.) A thicket? (\textbf{\ru{чаща}}.) Moss? (\textbf{\ru{мох}}.) A cone? (\textbf{\ru{шишка}}.) A tent? (\textbf{\ru{палатка}}.) A sleeping bag? (\textbf{\ru{спальник}}.) A campfire? (\textbf{\ru{костёр}}.) To light? (\textbf{\ru{разжигать}}.) To put out? (\textbf{\ru{тушить}}.) A high tide? (\textbf{\ru{прилив}}.) A low tide? (\textbf{\ru{отлив}}.) A bay? (\textbf{\ru{бухта}}.) A seashell? (\textbf{\ru{ракушка}}.) To dive? (\textbf{\ru{нырять}}.) A forecast? (\textbf{\ru{прогноз}}.) A weather forecaster? (\textbf{\ru{синоптик}}.) Precipitation? (\textbf{\ru{осадки}}.) Overcast? (\textbf{\ru{пасмурный}}.) To clear up? (\textbf{\ru{проясняться}}.) A downpour? (\textbf{\ru{ливень}}.) To drizzle? (\textbf{\ru{моросить}}.) To get soaked? (\textbf{\ru{промокать}}.) A puddle? (\textbf{\ru{лужа}}.) To pour? (\textbf{\ru{лить}}.) A storm? (\textbf{\ru{буря}}.) A hurricane? (\textbf{\ru{ураган}}.) A gust? (\textbf{\ru{порыв}}.) To rumble? (\textbf{\ru{греметь}}.) To flash? (\textbf{\ru{сверкать}}.)
\end{tcolorbox}
